
\subsection{模 $T^E$ 无障碍}
\label{subsec:6-5-5}
\liftIndexRangeStart{I376}

在定理 J 中，
我们假设 $L^{(0)},L^{(1)}$
经 bulk 形变后弱无障碍。
不过，在略弱的假设下也容易得到类似结论。

\begin{definition}{6.5.42}\label{def:6-5-42}
设 $L\subset M$ 是相对 spin
拉格朗日子流形。
若存在上链
$b\in C^1(L;\Lambda_{0,\mathrm{nov}}^+)$，
使
\begin{equation}
\widehat d(e^b)\equiv0
\pmod{T^E},
\tag{6.5.43.1}\label{eq:6-5-43-1}
\end{equation}
则称 $L$ \emph{模 $T^E$ 无障碍}\liftIndex{I320}。

若存在上链
$b\in C^1(L;\Lambda_{0,\mathrm{nov}}^+)^+$，
使
\begin{equation}
\widehat d(e^b)\equiv0
\pmod{(T^E,\mathbf e^+)},
\tag{6.5.43.2}\label{eq:6-5-43-2}
\end{equation}
则称 $L$ \emph{模 $T^E$ 弱无障碍}。

若存在
$b\in C^1(L;\Lambda_{0,\mathrm{nov}}^+)$
及
$\mathfrak b\in H^2(M;\Lambda_{0,\mathrm{nov}}^+)$，
使
\translatornote{源书实体页 404
在经 bulk 形变后的两项定义中
把 $\mathfrak b$ 写作属于
$H^2(L;\Lambda_{0,\mathrm{nov}}^+)$；
但它是环境 $M$ 的 bulk 类，
且紧接的引理--定义 6.5.44
明确使用 $H^2(M;\Lambda_{0,\mathrm{nov}}^+)$。
此处据同一参数的型别统一，
见疑误 A-0402。}
\begin{equation}
\mathfrak q(e^{\mathfrak b},e^b)
\equiv0\pmod{T^E},
\tag{6.5.43.3}\label{eq:6-5-43-3}
\end{equation}
则称 $L$ \emph{bulk 形变后模 $T^E$ 无障碍}。

若存在
$b\in C^1(L;\Lambda_{0,\mathrm{nov}}^+)^+$
及
$\mathfrak b\in H^2(M;\Lambda_{0,\mathrm{nov}}^+)$，
使
\begin{equation}
\mathfrak q(e^{\mathfrak b},e^b)
\equiv0\pmod{(T^E,\mathbf e^+)},
\tag{6.5.43.4}\label{eq:6-5-43-4}
\end{equation}
则称 $L$ \emph{bulk 形变后模 $T^E$ 弱无障碍}。
\end{definition}

满足式~\eqref{eq:6-5-43-2}
或~\eqref{eq:6-5-43-4} 时，
以
\[
\mathfrak{PO}(\mathfrak b,b)
\in
\Lambda_{0,\mathrm{nov}}/
T^E\Lambda_{0,\mathrm{nov}}
\]
记势函数，其定义关系为
\[
\mathfrak q(e^{\mathfrak b},e^b)
-\mathfrak{PO}(\mathfrak b,b)\mathbf e^+
\equiv0\pmod{T^E}.
\]
注意，对任意 $L$，
总存在某个正常数 $E$，
使 $L$ 模 $T^E$ 无障碍。

下一个引理的证明只需直接修改
命题~\ref{prop:3-7-17} 的证明。

\begin{statement}{引理--定义}{6.5.44}\label{lem-def:6-5-44}
假设 $L^{(0)}$ 与 $L^{(1)}$
干净相交。
若
$b_i\in C^1(L^{(i)};
\Lambda_{0,\mathrm{nov}}^+)$
及
$\mathfrak b\in H^2(M;
\Lambda_{0,\mathrm{nov}}^+)$
满足式~\eqref{eq:6-5-43-4}，
则定义
\[
\mathfrak m_1^{(\mathfrak b,b_1),(\mathfrak b,b_0)}:
C^*(L^{(1)},L^{(0)};
\Lambda_{0,\mathrm{nov}})
\longrightarrow
C^{*+1}(L^{(1)},L^{(0)};
\Lambda_{0,\mathrm{nov}})
\]
如下：
\[
\delta_{(\mathfrak b,b_1),(\mathfrak b,b_0)}(x)
=\mathfrak r(e^{\mathfrak b},e^{b_1}xe^{b_0}).
\]
若
$\mathfrak{PO}(\mathfrak b,b_0)
=\mathfrak{PO}(\mathfrak b,b_1)$，
则
\[
\delta_{(\mathfrak b,b_1),(\mathfrak b,b_0)}
\circ
\delta_{(\mathfrak b,b_1),(\mathfrak b,b_0)}
\equiv0\pmod{T^E}.
\]
\end{statement}

\begin{definition}{6.5.45}\label{def:6-5-45}
在引理--定义~\ref{lem-def:6-5-44}
的相同条件下，
定义 $\Lambda_{0,\mathrm{nov}}/
T^E\Lambda_{0,\mathrm{nov}}$-模
\[
HF\bigl((L^{(1)},(\mathfrak b,b_1)),
(L^{(0)},(\mathfrak b,b_0));
\Lambda_{0,\mathrm{nov}}/
T^E\Lambda_{0,\mathrm{nov}}\bigr)
\]
为商复形上所诱导算子的
\[
\frac{ker
\delta_{(\mathfrak b,b_1),(\mathfrak b,b_0)}}
{\operatorname{Im}
\delta_{(\mathfrak b,b_1),(\mathfrak b,b_0)}}.
\]
\end{definition}

仿照命题~\ref{prop:6-3-14} 的证明，
可证存在 $0<\lambda_i\le E$，使
\begin{equation}
\begin{aligned}
&HF\bigl((L^{(1)},(\mathfrak b,b_1)),
(L^{(0)},(\mathfrak b,b_0));
\Lambda_{0,\mathrm{nov}}/
T^E\Lambda_{0,\mathrm{nov}}\bigr)\\
&\qquad\cong
\bigoplus_i
\frac{\Lambda_{0,\mathrm{nov}}}
{T^{\lambda_i}\Lambda_{0,\mathrm{nov}}}.
\end{aligned}
\tag{6.5.46}\label{eq:6-5-46}
\end{equation}
\newpage
在这种情形，
定理 J 有如下推广。

\begin{theorem}{6.5.47}\label{thm:6-5-47}
在定义~\ref{def:6-5-45} 的假设下，
以式~\eqref{eq:6-5-46} 定义
$\lambda_i$，并置
\translatornote{源书实体页 405
将第二个计数函数写成 $a(i)$，
而结论使用 $a(\mu)$；
这里统一为后者，
见疑误 A-0403。}
\[
b(\mu)
=\#\{i\mid E>\lambda_i\ge\mu\},
\qquad
a(\mu)
=\#\{i\mid E=\lambda_i\}.
\]
若 $\mu$ 是哈密顿微分同胚
$\psi:M\to M$ 到恒等映射的
霍费尔距离，则
\translatornote{原书未限制 $\mu<E$，
但本节只掌握模 $T^E$ 的低能信息，
不能在能量损失达到 $E$ 后
继续保留 $\lambda_i=E$ 的顶端项；
无条件的下界可与可位移
拉格朗日子的例子冲突。
此处忠实保留源陈述，
其能量范围缺口见疑误 A-0404。}
\[
\#\bigl(\psi(L^{(0)})\cap L^{(1)}\bigr)
\ge a(\mu)+2b(\mu).
\]
\end{theorem}

原书称其证明只是对定理 J 的证明
作小幅修改，故予省略；
定理~\ref{thm:6-1-25}
也可以同样推广。
定理~\ref{thm:6-5-47}
略微推广了切卡诺夫在
\textbf{[Chek96,98]} 中的定理。
\liftIndexRangeEnd{I376}
